% Part: counterfactuals
% Chapter: introduction
% Section: material-conditional

\documentclass[../../../include/open-logic-section]{subfiles}

\begin{document}

\olfileid{cnt}{int}{mat}

\olsection{شرطی مادی}

شرطی در ساده‌ترین صورت انگلیسی خود جمله‌ای به شکل «اگر \dots، آنگاه
\dots» است که در آن جای خالی‌ها خود جمله‌اند؛ مانند «اگر پیشخدمت این
کار را کرده باشد، آنگاه باغبان بی‌گناه است». در درس‌های مقدماتی منطق
می‌آموزیم شرطی‌ها را با رابط~$\lif$ نمادگذاری کنیم: بخش‌های نشان‌داده‌شده
با~\dots{} را، مثلاً با !!{formula}های~$!A$ و~$!B$، نمادگذاری می‌کنیم
و کل شرطی با~$!A \lif !B$ نمادگذاری می‌شود.

رابط~$\lif$ \emph{تابعِ صدق} است؛ یعنی ارزش صدق---$\True$ یا
$\False$---ِ~$!A \lif !B$ با ارزش‌های صدق~$!A$ و~$!B$ تعیین می‌شود:
$!A \lif !B$ صادق است اگر و تنها اگر~$!A$ کاذب یا~$!B$ صادق باشد، و
در غیر این صورت کاذب است. نسبت به یک تخصیص ارزش صدق~$\pAssign v$،
$\pSat{v}{!A \lif !B}$ را اگر و تنها اگر~$\pSat/{v}{!A}$ یا
$\pSat{v}{!B}$ تعریف می‌کنیم. رابط~$\lif$ با این معناشناسی
\emph{شرطی مادی} نامیده می‌شود.

این تعریف شماری واقعیت منطقی ابتدایی به دست می‌دهد. نخست، قضیهٔ استنتاج
برای شرطی مادی برقرار است:
\begin{align}
  \text{اگر } \Gamma, !A \Entails !B \text{، آنگاه } \Gamma \Entails !A \lif !B
\end{align}
این شرطی تابعِ صدق است: $!A \lif !B$ و~$\lnot !A \lor !B$ هم‌ارزند:
\begin{align}
  !A \lif !B & \Entails \lnot !A \lor !B\\
  \lnot !A \lor !B & \Entails !A \lif !B
  \intertext{تالیِ شرطی مادی و نیز نقیض مقدم آن، شرطی را نتیجه می‌دهند:}
  !B & \Entails !A \lif !B\\
  \lnot !A & \Entails !A \lif !B
  \intertext{نقیض یک شرطی مادی با عطف مقدم آن و نقیض تالی آن هم‌ارز
    است: اگر~$!A \lif !B$ کاذب باشد، $!A \land \lnot !B$ صادق است و
    برعکس:}
  \lnot(!A \lif !B) & \Entails !A \land \lnot !B\\
  !A \land \lnot !B & \Entails \lnot(!A \lif !B)
  \intertext{شرطی مادی وضع مقدم را پشتیبانی می‌کند:}
  !A, !A \lif !B & \Entails !B
  \intertext{شرطی مادی تجمیع‌پذیر است:}
  !A \lif !B,  !A \lif !C & \Entails !A \lif (!B \land !C)
  \intertext{همواره می‌توان مقدم را تقویت کرد؛ یعنی شرطی \emph{یکنوا} است:}
  !A \lif !B & \Entails (!A \land !C) \lif !B
  \intertext{شرطی مادی تعدی است؛ یعنی قاعدهٔ زنجیره معتبر است:}
  !A \lif !B, !B \lif !C & \Entails !A \lif !C
  \intertext{شرطی مادی با عکس نقیض خود هم‌ارز است:}
  !A \lif !B & \Entails \lnot !B \lif \lnot !A\\
  \lnot !B \lif \lnot !A & \Entails !A \lif !B
\end{align}

این‌ها همگی استنتاج‌هایی سودمند و بی‌اشکال در استدلال ریاضی‌اند. اما
آثار فلسفی و زبان‌شناختی سرشار از پادنمونه‌های ادعایی برای استنتاج‌های
متناظر در بافت‌های غیرریاضی است. این پادنمونه‌ها نشان می‌دهند شرطی
مادی~$\lif$ رابط مناسب---یا دست‌کم همواره رابط مناسب---برای نمادگذاری
عبارت‌های انگلیسی «اگر \dots، آنگاه \dots» نیست.

\end{document}
